\section{Hypothesis, risks, and paper identity}
\label{sec:risks}

\subsection{Primary hypothesis}

\begin{quote}
When a probabilistic LLM judge captures useful but systematically imperfect information about human
preferences, the judge-to-human discrepancy can be statistically easier to estimate than human
preference from scarce labels alone. A minimax-adaptive distributional personalization estimator
should then reduce the human supervision needed to build an effective preference target, and this
statistical gain should carry over to better post-trained policies.
\end{quote}

The hypothesis can fail in informative ways. If the judge distribution carries nothing beyond a
scalar score ($G=0$), \dfsp{} should not materially beat \sfsp{} (E2). If the judge is misleading,
adaptive borrowing should revert to human-only estimation (E4). Both outcomes are reportable.

\subsection{Risks and mitigations}

\begin{center}
\small
\renewcommand{\arraystretch}{1.2}
\begin{tabularx}{\linewidth}{@{}p{4.4cm}X@{}}
\toprule
Risk & Mitigation \\
\midrule
Theory looks like a coordinatewise corollary of \citet{li2026personalizing} &
Rest the paper on the E2 finding and the acquisition result (Target~\ref{tr:design}); present
Targets~\ref{tr:lower} and~\ref{tr:e2e} and the propositions as supporting theory
(Remark~\ref{rem:novelty-theory}). \\
The judge protocol adds nothing beyond $p_J$ ($G=0$) &
Graded, two-order protocol J1; test $G$ with calibrated variable-importance inference (E2); gate G1. \\
Targets stated without necessary conditions &
Review counterexamples are recorded next to Targets~\ref{tr:upper} and~\ref{tr:lower}; keep the
nondegeneracy, occupancy, support, compatibility and slack conditions explicit. \\
Diagnostic or test labels leak into budgeted procedures &
Separate diagnostic ($\mathcal G$), development ($\mathcal V$) and test ($\mathcal T$) splits; freeze
procedures after the label-free feasibility pilot. \\
Curse of dimensionality in $Z$ &
Localize on low-dimensional judge summaries ($d\le4$); these are natural covariates (Design D1). \\
Judge support differs from the human label space &
Use the link $T$; the distribution then enters through the localization channel
(Section~\ref{sec:two-channels}). \\
Test contamination across arms &
Common prompt-level test set $\mathcal T$ fixed before acquisition and training. \\
Single-label data cannot evaluate distribution-level metrics &
Lemma~\ref{lem:heldout} makes Brier and RPS valid for comparing methods; use distribution-level
metrics only on multi-annotator subsets. \\
Preference-loss gains do not become generation gains &
Report P2 and P3 alongside P1; state the scope of the bridge explicitly. \\
Evaluation leakage (judge self-preference) &
Gold reward model trained on human labels; an evaluation judge different from the training judge. \\
Fast-moving neighbors (e.g.\ \citealp{cai2026dial}, posted 25 Sep 2026) &
Re-run the novelty check before submission. Differentiate on per-comparison, nonparametric,
distribution-valued targets with a post-training bridge. \\
\bottomrule
\end{tabularx}
\end{center}

\subsection{Expected contributions}

\begin{enumerate}[label=\textbf{C\arabic*.}]
  \item \emph{Distributional few-shot personalization.} Personalization of a graded,
        two-order judge distribution with an estimator that adaptively borrows from a frozen judge
        using a few human labels, together with a calibrated test of whether the distribution adds
        information beyond $p_J$.
  \item \emph{Label design under the population risk.} A finite-pool acquisition rule driven by pool
        density, human variability and local discrepancy complexity, with a conditional
        upper-bound guarantee. It improves constants, not rates.
  \item \emph{Post-training evidence and supporting theory.} Real checkpoints trained against the
        personalized target, with policy quality reported against the human-label budget. The
        supporting theory covers bounded-target minimax lower bounds with a fixed judge, the bridges,
        and their end-to-end corollary.
\end{enumerate}

\subsection{Working titles and message}

\begin{itemize}
  \item \emph{Distributional Few-Shot Personalization for Human-Anchored LLM Post-Training}
  \item \emph{From AI Distributions to Human Preferences: Minimax Few-Shot Personalization for LLM
        Post-Training}
\end{itemize}

\keybox{\centering Do not treat scarce human feedback merely as corrections to AI labels.
Use it to statistically personalize the entire AI preference model, then post-train the policy
against that personalized human target.}
